\documentclass[11pt]{article}

% --- Preámbulo compatible con arXiv (solo paquetes estándar) ---
\usepackage[utf8]{inputenc}
\usepackage[T1]{fontenc}
\usepackage[spanish]{babel}
\usepackage{amsmath,amssymb,amsthm}
\usepackage{booktabs}
\usepackage{geometry}
\usepackage{hyperref}
\usepackage{xcolor}
\geometry{margin=1in}
\hypersetup{colorlinks=true,linkcolor=blue,citecolor=blue,urlcolor=blue}

\newtheorem{definition}{Definición}
\newtheorem{proposition}{Proposición}
\newtheorem{theorem}{Teorema}

\newcommand{\dcausal}{d_{\mathrm{c}}}
\newcommand{\Nk}{\mathcal{N}_k}
\newcommand{\code}[1]{\texttt{#1}}

\title{\textbf{La MMU cognitiva:\\ memoria de trabajo determinista y auditable para agentes de código basados en LLM\\ \large (y un relato honesto de por qué la recuperación causal no supera a RAG)}}
\author{El proyecto CCOS\thanks{Causal Context Operating System (CCOS), un prototipo
de investigación abierto. Código fuente, guiones de reproducción y este artículo:
\url{https://github.com/CHECKUPAUTO/CCOS}.}}
\date{\today}

\begin{document}
\maketitle

\begin{abstract}
La memoria de trabajo de un agente de código basado en LLM se gestiona de forma
opaca: una pila de recuperación (RAG, GraphRAG, marcos de agentes) pagina un
repositorio en una ventana acotada y muta esa ventana de manera probabilística a lo
largo de una sesión, de modo que cuando un agente se desvía tras decenas de pasos es
imposible decir \emph{por qué} ni \emph{cuándo} se corrompió su representación del
proyecto. CCOS trata la memoria del agente como un sistema operativo trata la RAM: el
código fuente se analiza en un grafo causal \emph{event-sourced}, los nodos se
puntúan y se paginan bajo un presupuesto de tokens, y cada transición de estado se
registra en un registro encadenado por hash que se reproduce bit a bit. Damos una
definición precisa y determinista de \emph{región causal} (una componente conexa del
grafo de enlaces causales entre archivos) y de distancia causal (un camino más corto
ponderado), y probamos que las regiones y la ventana paginada son una función pura
del grafo, reconstruyéndose exactamente en la reproducción --- la propiedad que hace
que la memoria de un agente sea \emph{auditable} y \emph{depurable por viaje en el
tiempo}. Luego preguntamos si esta estructura causal hace de CCOS un \emph{mejor
recuperador} que una referencia básica, e informamos --- honestamente --- de que no.
Medida sobre $70$ commits reales de corrección de errores en tres bibliotecas
(«crates») maduras, la selección causal \textbf{empata con un simple recuperador
léxico TF-IDF} y pierde ante él con un presupuesto ajustado; en errores reales los
archivos de una corrección comparten vocabulario, así que la similitud léxica también
los encuentra, y una prueba de concepto que arranca CCOS desde una traza de fallo es
asimismo superada por un RAG-sobre-el-mensaje-de-error (los mensajes de error de Rust
nombran la causa). Una simulación sin LLM \emph{sí} separa a ambos cuando una causa se
construye para ser léxicamente distinta de su síntoma, y una primera ejecución en
dispositivo con un modelo de $7$ mil millones de parámetros lo reproduce --- pero ese
régimen no ocurre en código real. La contribución no es, por tanto, una mejor
recuperación. Es la infraestructura que a una pila de recuperación le falta
estructuralmente: una memoria de agente \emph{determinista, reproducible, auditable}
en la que se puede rebobinar el estado exacto del contexto en cualquier paso y
reproducirlo bajo parámetros distintos (\emph{depuración por viaje en el tiempo}).
Hasta donde sabemos, CCOS es el \emph{primer} sistema que trata la \emph{propia
memoria de trabajo} de un agente como un subsistema transaccional --- determinista,
encadenado por hash, reproducible bit a bit y \emph{depurable a posteriori}: una
caja negra de la atención, con un punto de interrupción en el paso exacto en que la
información correcta fue desalojada de la ventana. Publicamos el motor, el arnés de validación honesto y este artículo, etiquetando en
todo momento lo que está \emph{probado} (estructura formal, determinismo,
reproducción), \emph{medido} (cobertura de recuperación --- donde CCOS \emph{no}
supera al RAG léxico en errores reales), y el único eje en el que CCOS \emph{sí}
gana: la \textbf{eficiencia}. Ejecutado de extremo a extremo sobre correcciones
reales de un solo archivo con un modelo de 30B y un reintento con el compilador en el
bucle, CCOS y RAG resuelven \emph{por igual} ($2/10$), pero CCOS lo hace con $776$
tokens de contexto frente a los $5366$ del recuperador léxico --- una reducción de
$6{,}9\times$ (y de $4{,}1$ a $9{,}1\times$ en $51$ correcciones de tres crates, sin
modelo), porque la región causal se detiene en el working set en lugar de rellenar un
top-$k$ hasta el presupuesto. CCOS no es, pues, un mejor recuperador ni solucionador,
sino uno más \emph{frugal} y \emph{auditable} --- se autocalibra, sin $k$ que ajustar.
(Sobre su propio árbol de fuentes, la selección por región cubre el vecindario a $k$
saltos de una tarea con una exhaustividad de $0{,}97$ frente a $0{,}35$ de la
selección plana, con $\approx\!48\%$ menos tokens; las regiones están conectadas
internamente al $95{,}5\%$.)
\end{abstract}

\section{Introducción}

Un agente de código basado en LLM que opera sobre un repositorio debe decidir
repetidamente \emph{qué poner en la ventana de contexto}. El patrón dominante es la
recuperación: incrustar fragmentos de código, clasificarlos frente a la tarea, y
concatenar el top-$k$ hasta agotar un presupuesto de tokens~\cite{lewis2020rag}. Esto
trata el contexto como una lista clasificada unidimensional. De ahí se siguen dos
modos de fallo. Primero, la \emph{dilución}: a medida que las tareas se alargan, la
ventana acumula código globalmente saliente pero irrelevante para la tarea, y los
modelos atienden mal al medio de los contextos largos~\cite{liu2023lost}. Segundo, la
\emph{incoherencia}: los $k$ fragmentos mejor clasificados no tienen por qué formar
una unidad conexa de razonamiento --- una función puede paginarse sin sus llamadores,
sus sitios de error o los datos de los que depende.

Los sistemas operativos afrontaron el problema análogo hace décadas y respondieron
con el \emph{working set}: el conjunto de páginas referenciadas recientemente que
deben permanecer residentes~\cite{denning1968working}. CCOS adopta la analogía para
el contexto de LLM~\cite{packer2023memgpt}: el código fuente se analiza en un grafo
causal, los nodos se puntúan, y una «ventana de contexto» acotada se pagina dentro y
fuera como RAM~$\leftrightarrow$~VRAM. A diferencia de una pila de recuperación, cada
transición se registra en un registro de eventos encadenado por hash, de modo que la
memoria no es una caja negra sino un artefacto \emph{auditable y reproducible}.
Nuestras contribuciones son:

\begin{enumerate}
  \item Un \textbf{modelo formal y determinista de regiones de contexto}
  (\S\ref{sec:regions}, \S\ref{sec:determinism}): la distancia causal como camino más
  corto ponderado, la pertenencia a una región como componente conexa del grafo de
  enlaces causales entre archivos, y un teorema de determinismo --- las regiones y la
  ventana paginada son una función pura del grafo, así que una sesión se reconstruye
  bit a bit desde su registro de eventos encadenado por hash (reproducción a prueba de
  manipulación).
  \item \textbf{Sesiones de agente event-sourced con depuración por viaje en el
  tiempo} (\S\ref{sec:timetravel}): cada operación cognitiva (ingesta, señal de fallo,
  recuerdo) queda registrada; el estado exacto del contexto en cualquier paso es
  reconstruible; y un recuerdo puede \emph{reproducirse bajo parámetros distintos}
  para preguntar si el agente habría decidido mejor --- la capacidad que a una pila de
  recuperación probabilística le falta estructuralmente. Hasta donde sabemos, es el
  primer tratamiento del \emph{ensamblaje del contexto en sí} como un subsistema
  reproducible y depurable a posteriori --- una \emph{caja negra de la atención de un
  agente} --- con un \emph{punto de interrupción de desalojo} que nombra el paso y la
  operación exactos en que la causa verdadera fue expulsada de la ventana acotada.
  \item Un \textbf{arnés de validación honesto y un resultado negativo}
  (\S\ref{sec:protocol}): sobre $70$ commits reales de corrección, la selección causal
  \emph{no} supera a un recuperador léxico TF-IDF al colocar los archivos de una
  corrección en la ventana (empata, y pierde con un presupuesto ajustado), y un pivote
  por traza de fallo es superado por un RAG-sobre-el-mensaje-de-error. Lo informamos
  con claridad --- reubica el valor de CCOS de la \emph{recuperación} a la
  \emph{auditabilidad}.
  \item Una \textbf{medición de localidad sin LLM} (\S\ref{sec:eval}) con cifras
  reales y reproducibles, y un protocolo falsable de condición \emph{suficiente}
  (Fase~4) que especificamos pero dejamos abierto.
\end{enumerate}

\section{Trabajo relacionado}\label{sec:related}

\paragraph{Generación aumentada por recuperación.} RAG aumenta un modelo paramétrico
con un almacén no paramétrico y recupera pasajes por consulta~\cite{lewis2020rag}.
Self-RAG añade tokens de reflexión que regulan la recuperación y critican las
generaciones~\cite{asai2023selfrag}. Estos operan sobre fragmentos
\emph{independientes}; la coherencia entre elementos recuperados no se modela.

\paragraph{Recuperación consciente de grafos y estructura.} GraphRAG construye un
grafo de conocimiento de entidades y responde consultas globales resumiendo la
estructura comunitaria~\cite{edge2024graphrag}. Para código, los grafos de propiedades
unifican AST, flujo de control y flujo de datos~\cite{yamaguchi2014cpg}. Las regiones
de CCOS están en este linaje pero apuntan a la \emph{paginación}: qué subestructura
conexa hacer residente para una tarea, bajo un presupuesto de tokens, con desalojo
determinista.

\paragraph{Memoria de agente.} MemGPT presenta el LLM como un SO que pagina entre una
ventana en contexto y almacenamiento externo~\cite{packer2023memgpt}; los Generative
Agents recuperan recuerdos puntuados por recencia, importancia y
relevancia~\cite{park2023generative} --- los mismos factores que CCOS agrega en una
temperatura de región. LangGraph~\cite{langgraph} estructura los agentes como grafos
de pasos con estado; orquesta el flujo de \emph{control}, mientras que CCOS estructura
la \emph{memoria}. Ambos son complementarios.

\paragraph{Gestión de la ventana de contexto.} El desalojo de la caché KV mantiene
residentes los tokens «grandes impactadores» o sumidero~\cite{liu2023lost}, y
PagedAttention aplica la paginación de SO a la caché KV~\cite{kwon2023paged}. Estas
técnicas actúan a nivel de token dentro de un único paso hacia adelante; CCOS actúa al
nivel \emph{semántico} a lo largo de una sesión de agente.

\section{El sustrato de contexto causal}\label{sec:substrate}

CCOS analiza el código fuente en un \emph{grafo de memoria causal} dirigido $G=(V,E)$.
Un nodo $v\in V$ es un archivo, módulo, símbolo o dependencia externa, con campos
escalares $\mathrm{imp}(v)$ (importancia base), $\mathrm{fail}(v)\in[0,1]$ (relevancia
de fallo) y $\mathrm{rec}(v)\in[0,1]$ (recencia). Una arista $e=(u\!\to\!w)\in E$ lleva
un peso $w(e)\in(0,1]$ y un tipo (contención, dependencia, referencia, causación). El
núcleo asigna a cada nodo una puntuación causal
\begin{equation}
\mathrm{score}(v) = \mathrm{clamp}\big(0.15\,\mathrm{imp}(v) + 0.50\,\mathrm{fail}(v)
  + 0.30\,\mathrm{rec}(v) + 0.05\ln(1{+}\mathrm{acc}(v)),\,0,1\big),
\label{eq:score}
\end{equation}
donde $\mathrm{acc}(v)$ es el conteo de accesos. Los fallos se propagan a lo largo de
las aristas, $\mathrm{fail}$ decae con un reloj lógico, y cada transición de estado se
añade a un registro de eventos encadenado por hash que permite una reproducción
determinista~\cite{haber1991timestamp}. Los identificadores de nodo tienen espacio de
nombres (\code{file:p}, \code{mod:p:n}, \code{use:p:path}, \code{sym:p:n},
\code{dep:root}); el archivo propietario de un nodo es recuperable desde su
identificador. CCOS v0.2 pagina los nodos de mayor $\mathrm{score}$: una política
plana y 1-D. Ahora hacemos espacial la selección.

\section{Regiones de contexto}\label{sec:regions}

\subsection{Distancia causal}

\begin{definition}[Distancia causal]
Sea $\hat G$ el multigrafo no dirigido sobre $V$ inducido por $E$, y asignemos a cada
arista el coste $c(e) = -\ln w(e) \ge 0$, de modo que un enlace causal más fuerte sea
un paso más corto. La \emph{distancia causal} $\dcausal(u,v)$ es el coste total mínimo
sobre todos los caminos $u$--$v$ en $\hat G$, y $+\infty$ si no existe ninguno. La
distancia en saltos no ponderada $\mathrm{hops}(u,v)$ se define análogamente con
costes unitarios.
\end{definition}

$c(e)\ge 0$ puesto que $w(e)\le 1$, así que $\dcausal$ es una verdadera métrica de
camino más corto en cada componente conexa (no negativa, simétrica, desigualdad
triangular). El \emph{vecindario causal a $k$ saltos} de un objetivo $t$ es
\begin{equation}
\Nk(t) = \{\, v \in V : \mathrm{hops}(t,v) \le k \,\}.
\end{equation}
$\Nk(t)$ es la verdad-base «lo que es causalmente relevante para una tarea en $t$»
usada en la evaluación (\S\ref{sec:eval}).

\subsection{Pertenencia a una región}

Los concentradores de dependencias externas (p.\ ej.\ \code{dep:std}) conectan casi
todo y no deben colapsar el grafo en una sola componente. Por eso los separamos.

\begin{definition}[Enlace causal entre archivos]
Para un nodo no externo $v$ sea $\phi(v)$ su archivo propietario. Dos archivos $f,g$
están \emph{directamente enlazados}, escrito $f \approx g$, si y solo si existe una
arista $(u\!\to\!w)\in E$ con $\phi(u)=f$, $\phi(w)=g$, $f\neq g$, y ni $u$ ni $w$ es
un nodo de dependencia externa.
\end{definition}

\begin{definition}[Región]\label{def:region}
Sea $\approx^{*}$ la clausura reflexiva-transitiva de $\approx$ sobre el conjunto de
claves de archivo. Una \emph{región} es el conjunto de todos los nodos cuyos archivos
están en una misma clase de equivalencia de $\approx^{*}$; todos los nodos de
dependencia externa forman una región adicional. Dos nodos pertenecen a la misma
región si y solo si sus archivos están en la misma componente conexa del grafo de
enlaces causales entre archivos.
\end{definition}

\begin{proposition}[Las regiones particionan los nodos]
Las regiones de la Definición~\ref{def:region} forman una partición de $V$: cada nodo
está en exactamente una región.
\end{proposition}
\begin{proof}
$\approx^{*}$ es una relación de equivalencia sobre las claves de archivo, así que sus
clases particionan las claves; el cubo externo es una clase distinta por construcción.
Aplicar cada nodo no externo a la clase de su archivo y cada nodo externo a la clase
externa es una función total hacia bloques disjuntos, y por tanto una partición. \qed
\end{proof}

Por defecto (solo aristas de contención e importación) cada archivo fuente es su
propia región. Una verdadera dependencia entre archivos o un fallo propagado fusiona
archivos en una única región multiarchivo --- la «zona de conocimiento» que un agente
debería despertar junta.

\subsection{Escalares de región}

Para una región $R$ con conjunto de miembros $M$:
\begin{align}
\mathrm{heat}(v)        &= 0.5\,\mathrm{score}(v) + 0.3\,\mathrm{fail}(v) + 0.2\,\mathrm{rec}(v), \\
\mathrm{temp}(R)        &= \mathrm{clamp}\!\Big(\tfrac{1}{|M|}\textstyle\sum_{v\in M}\mathrm{heat}(v),\,0,1\Big), \label{eq:temp}\\
\mathrm{dens}(R)        &= \frac{|\{\,e\in E : \text{ambos extremos} \in M\,\}|}{|M|}. \label{eq:dens}
\end{align}
La temperatura es cuán «despierta» está una región; la densidad es su cohesión causal
interna (aristas internas por miembro). La activación calienta una región
($\mathrm{temp}\mathrel{+}=0.25$, con tope) y registra un tic lógico; un paso de
enfriamiento multiplica las temperaturas por un factor de decaimiento y desaloja las
regiones por debajo de un umbral.

\subsection{Política de admisión dinámica}\label{sec:policy}

El umbral estático $0.6$ se convierte en una función de la presión de tokens
$u\in[0,1]$ (la fracción usada del presupuesto) y la complejidad de la tarea
$\kappa\in[0,1]$:
\begin{align}
\theta(u,\kappa)       &= \mathrm{clamp}(0.6 + 0.3\,u - 0.2\,\kappa,\; 0.05,\; 0.95), \\
a(R)                   &= 0.55\,\mathrm{temp}(R) + 0.30\,\frac{\mathrm{dens}(R)}{1+\mathrm{dens}(R)} + 0.15\,\kappa, \\
\mathrm{admit}(R)      &\iff a(R) \ge \theta(u,\kappa).
\end{align}
Una región caliente y cohesiva puede admitirse aun cuando el $0.6$ estático la
rechazaría; una ventana casi llena eleva $\theta$ para que solo entren las regiones
más calientes.

\section{Determinismo y reproducción}\label{sec:determinism}

\begin{theorem}[Determinismo regional]
La partición en regiones, cada escalar de región de las
ecuaciones~\eqref{eq:temp}--\eqref{eq:dens}, y el historial de activación son
funciones puras, independientes del orden, del grafo $G$ y de la secuencia de eventos
de activación (lógicamente marcados con tiempo). En consecuencia, dado el registro de
eventos de una sesión, el estado del motor se reconstruye idénticamente: si $G'$ es el
grafo reconstruido desde el registro y $L$ sus eventos de región, entonces
$\mathrm{replay\_from}(G',L)$ es igual al motor en vivo que produjo $L$.
\end{theorem}
\begin{proof}[Esbozo de prueba]
El agrupamiento enumera nodos y aristas en orden ordenado y calcula las componentes
conexas mediante una búsqueda en anchura ordenada, de modo que su salida es
independiente del orden de iteración del hash. Las ecuaciones \eqref{eq:score},
\eqref{eq:temp} y \eqref{eq:dens} son aritmética determinista sobre los campos de los
nodos y un conteo de aristas que cualifican. La activación lee un reloj lógico (un
contador), nunca el tiempo de pared, y el evento \code{RegionActivated} emitido
registra el tic exacto. La reconstrucción reagrupa $G'$ (estado base idéntico, ya que
$G'\!=\!G$ estructuralmente) y aplica las activaciones registradas en el orden del
registro; cada paso es la misma función pura de las mismas entradas, así que el estado
resultante es idéntico. La prueba de integración
\code{replay\_reconstructs\_identical\_engine} verifica
$\mathrm{engine}=\mathrm{replay\_from}(G',L)$ por igualdad estructural. \qed
\end{proof}

Esto extiende las garantías existentes de CCOS: el registro de eventos principal está
encadenado por hash y es a prueba de manipulación, de modo que el historial de
regiones es auditable, y una prueba de 10\,000 ciclos no exhibe deriva alguna en el
número de regiones ni en las temperaturas.

\paragraph{Una frontera de entrada auditable.} El mismo sustrato determinista y
reproducible endurece el texto que un agente ingiere. Los vectores de inyección
por caracteres ocultos --- anulaciones bidireccionales (el ataque \emph{Trojan
Source}, CVE-2021-42574), formato de ancho cero y el bloque Unicode \emph{Tags}
usado para el contrabando de ASCII invisible --- se des-ofuscan en la ingesta a
literales explícitos y visibles (\code{[U+202E RLO]}) en lugar de eliminarse en
silencio, y los hallazgos se registran en el mismo registro encadenado por hash,
de modo que una reproducción reconstruye el estado limpio. Un clasificador lineal
en espacio logarítmico posterior --- la forma cerrada del Naïve Bayes multinomial,
$\mathrm{logit}=b+W\!\cdot\!X$ sobre un vector de características por \emph{hashing
trick}, con sus pesos fijados en un blob verificado por suma de comprobación ---
añade una señal de inyección \emph{determinista y descomponible forensemente} (un
red-team sobre datos reservados mide $F_1=0{,}90$; precisión $0{,}87$,
exhaustividad $0{,}93$). Lo afirmamos de forma deliberada: es una \emph{señal, no
una solución} --- un modelo de bolsa de características se evade con paráfrasis, y
ninguna pasada a nivel de carácter aborda la inyección semántica; la mitigación
real es la separación de privilegios en el host. La contribución no es un detector
novedoso, sino que la des-ofuscación y la puntuación heredan la propiedad
distintiva de CCOS: toda decisión de seguridad es reproducible y auditable hasta
la característica exacta que la inclinó.

\section{Sesiones event-sourced y depuración por viaje en el tiempo}\label{sec:timetravel}

El determinismo no es solo una propiedad de corrección; es la base de la capacidad
que, como argumentará nuestra evaluación, distingue a CCOS de una pila de
recuperación. Una \code{AgentSession} registra cada operación cognitiva que un agente
realiza sobre su memoria --- \code{Ingest}, \code{SignalFailure}, \code{Recall} ---
como una línea de tiempo ordenada. Dado que cada operación es una función determinista
del estado previo, la memoria tras las primeras $n$ operaciones se reconstruye
exactamente reproduciendo las que mutan (\code{replay\_to(n)}): se puede
\emph{rebobinar la mente del agente} a cualquier paso.

La operación que importa para la depuración es la contrafactual.
\code{recall\_what\_if($n$, $q$, $b$)} reproduce la memoria hasta el paso $n$ y vuelve
a ejecutar un recuerdo con una consulta $q$ o un presupuesto de tokens $b$ distintos,
devolviendo la ventana que el agente \emph{habría} visto. Cuando un agente emite un
mal parche en el paso 15, un operador puede reproducir su contexto exacto en el paso
14, ampliar el presupuesto o cambiar el ancla, y observar si la decisión mejora --- un
depurador de bucle cerrado y reproducible para la memoria de trabajo de un agente.

\paragraph{Un punto de interrupción sobre la atención.} La reproducción también
localiza la \emph{deriva}. El punto de observación \code{missing} recorre la línea de
tiempo en busca del paso exacto en que un nodo dado --- típicamente la causa
verdadera de un fallo --- fue expulsado de la ventana acotada por la presión
competidora, nombrando la operación desencadenante y la diferencia de tokens; una
vista \code{energy} complementaria revela la migración de calor causal por el grafo
que un diff a nivel de archivo pasa por alto. Es, en efecto, un \emph{punto de
interrupción sobre la atención de un agente}: el instante preciso en que la
información correcta abandonó la ventana. Una pila de recuperación opaca solo puede
mostrar el contexto final, corrupto; CCOS puede señalar el paso --- y la operación ---
en que se torció.

Una pila RAG o de marco muta su almacén de forma probabilística a lo largo de una sesión y
no conserva ningún registro de transiciones canónico y reproducible, así que la misma
pregunta (\emph{¿por qué, y cuándo, se corrompió la representación?}) no tiene
respuesta allí. No afirmamos que esto mejore el éxito en las tareas; afirmamos que es
una propiedad que las alternativas no tienen, y que es el lugar honesto de la
contribución de CCOS (\S\ref{sec:protocol}).

\section{Implementación}\label{sec:impl}

El motor son $\approx\!600$ líneas de Rust ligero en dependencias, en capa sobre el
grafo, sin cambios en el analizador, el guardián, el constructor incremental o el
núcleo event-sourcing. La superficie pública es: \code{ContextRegionEngine}
(\code{cluster\_nodes}, \code{initialize\_regions}, \code{activate\_region},
\code{tick\_cooldown}, \code{replay\_from}), los tipos de datos \code{ContextRegion} /
\code{ContextPoint} / \code{ContextWindow}, la \code{ContextPolicy}, y cinco nuevas
variantes de evento (\code{RegionCreated}, \code{RegionActivated}, \code{RegionMerged},
\code{RegionEvicted}, \code{ContextWindowGenerated}). Una CLI \code{ccos regions}
expone el agrupamiento, la activación y el informe de localidad. Toda la crate compila
sin advertencias bajo \code{clippy -D warnings} y pasa 364 pruebas.

\section{Evaluación: lo que podemos medir hoy}\label{sec:eval}

Separamos los resultados \emph{medidos} (esta sección, totalmente reproducible y sin
LLM) de las ganancias \emph{hipotéticas} a nivel de agente (\S\ref{sec:protocol}).

\paragraph{Montaje.} Operamos sobre el propio árbol de fuentes de CCOS ($|V|=705$
nodos, $|E|=822$ aristas, dando $26$ regiones). Para un nodo objetivo $t$ tomamos
$\Nk(t)$ con $k=2$ como verdad-base y comparamos dos estrategias de selección al
presupuesto de tamaño de la región: \textbf{plana} --- los $|R|$ mejores nodos por
$\mathrm{score}$ (CCOS v0.2 / recuperación clasificada clásica) --- y \textbf{región}
--- los miembros de la región de $t$. Reportamos precisión causal $|S\cap\Nk|/|S|$,
exhaustividad $|S\cap\Nk|/|\Nk|$, y los tokens que cada estrategia necesita para
\emph{cubrir} $\Nk$. Todas las cifras las emite \code{scripts/region\_benchmark.sh} y
son deterministas.

\begin{table}[h]
\centering
\begin{tabular}{lcc}
\toprule
Estrategia & precisión causal & exhaustividad causal \\
\midrule
plana (v0.2 / clasificada) & 0.021 & 0.347 \\
\textbf{región (v0.3)} & \textbf{0.057} & \textbf{0.972} \\
\bottomrule
\end{tabular}
\caption{Localidad sobre \code{src/} ($k=2$, 12 objetivos). La selección por región
cubre el $97\%$ del vecindario causal de una tarea frente al $35\%$ de la selección
plana, con $\approx\!48\%$ menos tokens para igual cobertura.}
\label{tab:locality}
\end{table}

\paragraph{Cohesión y coste.} Las regiones alcanzan una \textbf{densidad causal media
de $0{,}955$} (Ec.~\ref{eq:dens}, normalizada) --- están casi por completo conectadas
internamente, es decir, agrupaciones causales genuinas y no grupos de archivos
arbitrarios. Construir el mapa de regiones para $705$ nodos lleva $\approx\!20$\,ms
($54{,}5$ construcciones/s); una sola activación son $\approx\!20$\,ms.

\paragraph{Lectura honesta.} La precisión absoluta es baja ($0{,}06$) porque el
analizador v0.2 solo emite aristas de contención e importación, así que $\Nk(t)$ es
diminuto (a menudo $2$--$3$ nodos) mientras que una región abarca un archivo entero.
Las ganancias robustas y reales son la \emph{exhaustividad} ($0{,}97$ frente a
$0{,}35$) y la \emph{eficiencia en tokens} ($\approx\!48\%$): una región contiene de
forma fiable el vecindario causal de una tarea y paga menos tokens por hacerlo.
Aristas semánticas más ricas (grafo de llamadas / flujo de datos, ítem P1.3 de la hoja
de ruta) afinarían $\Nk$ y son la palanca principal para elevar la precisión; lo
señalamos en vez de ocultarlo.

\section{Simulación de hipótesis bajo un oráculo declarado (sin LLM)}\label{sec:sim}

La tesis central --- \emph{¿ayuda la memoria causal por regiones a un agente en
tareas largas multiarchivo?} --- en última instancia necesita despliegues de LLM
(\S\ref{sec:protocol}). Pero su \emph{condición necesaria} es comprobable ya, sin LLM:
\textbf{un agente no puede resolver una tarea cuyo contexto causal requerido está
ausente de su ventana}. Medimos, bajo un oráculo explícito, si cada estrategia de
selección \emph{coloca el contexto causal requerido en la ventana}.

\paragraph{Montaje.} Generamos repositorios sintéticos modulares: $S$ subsistemas
independientes, cada uno un conjunto de archivos enlazados por una cadena causal entre
archivos más símbolos \emph{señuelo} de alta puntuación; no hay aristas entre
subsistemas, así que cada subsistema es una región causal acotada. La estructura causal
está \emph{desacoplada} de la estructura léxica --- un vecino de la cadena no comparte
ningún token de identificador con el objetivo, solo una arista --- modelando una
dependencia causalmente esencial pero léxicamente distinta. Una tarea de \emph{diámetro}
$d$ requiere la ventana $\pm d$ a lo largo de su cadena (hasta $2d{+}1$ archivos); el
presupuesto contiene $\approx$ un subsistema, mucho menos que el repositorio. El
\textbf{oráculo} es $R(t)\subseteq S$.

Crucialmente, los \emph{pipelines} de recuperación localizan código a partir de una
\emph{consulta} textual, mientras que una memoria de tipo SO se \emph{ancla} en la
señal del espacio de trabajo (el archivo activo, una prueba que falla). Modelamos
ambos: cada tarea lleva una consulta (una bolsa de tokens) y un ancla (el nodo del
archivo activo), y ejecutamos dos escenarios --- \textbf{limpio} (la consulta apunta
al objetivo) y \textbf{ruidoso} (un señuelo trampa en un subsistema no relacionado
supera léxicamente al objetivo). Seis estrategias comparten el presupuesto:
\code{rag-dense} (top-$k$ léxico), \code{rag-hybrid} (léxico $+$ puntuación causal),
\code{graphrag-1hop} (mejor acierto $+$ un salto), \code{graphrag-bfs} (expansión no
acotada desde el mejor acierto), \code{ccos-from-query} (región CCOS del mejor acierto
\emph{léxico} --- una ablación), y \code{ccos-region} (región CCOS del \emph{ancla}).
Con semilla y deterministas; reproducir con \code{ccos experiment}.

\begin{table}[h]
\centering
\begin{tabular}{lcccccc}
\toprule
& \multicolumn{4}{c}{éxito al diámetro $d$} & \multicolumn{2}{c}{global} \\
\cmidrule(lr){2-5}\cmidrule(lr){6-7}
Estrategia & $d{=}1$ & $d{=}2$ & $d{=}3$ & $d{=}4$ & éxito & cob. \\
\midrule
\multicolumn{7}{l}{\emph{Consulta limpia (apunta al objetivo)}}\\
\code{rag-dense}      & 0.00 & 0.00 & 0.00 & 0.00 & 0.00 & 0.19 \\
\code{rag-hybrid}     & 1.00 & 0.00 & 0.00 & 0.00 & 0.23 & 0.65 \\
\code{graphrag-1hop}  & 1.00 & 0.00 & 0.00 & 0.00 & 0.23 & 0.58 \\
\code{graphrag-bfs}   & 1.00 & 1.00 & 1.00 & 1.00 & 1.00 & 1.00 \\
\code{ccos-from-query}& 1.00 & 1.00 & 1.00 & 1.00 & 1.00 & 1.00 \\
\code{ccos-region}    & 1.00 & 1.00 & 1.00 & 1.00 & 1.00 & 1.00 \\
\midrule
\multicolumn{7}{l}{\emph{Consulta ruidosa (un señuelo supera léxicamente al objetivo)}}\\
\code{rag-dense}      & 0.00 & 0.00 & 0.00 & 0.00 & 0.00 & 0.19 \\
\code{rag-hybrid}     & 1.00 & 0.00 & 0.00 & 0.00 & 0.23 & 0.65 \\
\code{graphrag-1hop}  & 0.00 & 0.00 & 0.00 & 0.00 & 0.00 & 0.19 \\
\code{graphrag-bfs}   & 0.00 & 0.00 & 0.00 & 0.00 & 0.00 & 0.00 \\
\code{ccos-from-query}& 0.00 & 0.00 & 0.00 & 0.00 & 0.00 & 0.00 \\
\textbf{\code{ccos-region}} & \textbf{1.00} & \textbf{1.00} & \textbf{1.00} & \textbf{1.00} & \textbf{1.00} & \textbf{1.00} \\
\bottomrule
\end{tabular}
\caption{Simulación de hipótesis ($800$ tareas, semilla $42$, presupuesto $\approx$ un
subsistema). Éxito $=$ el conjunto causal requerido $R(t)$ está dentro de la ventana.
Bajo una consulta limpia, todo método consciente de la estructura empata; bajo una
consulta engañosa, solo sobrevive la región anclada al espacio de trabajo.}
\label{tab:sim}
\end{table}

\paragraph{Hallazgos (y sus límites honestos).}
\textbf{(1)} La recuperación léxica (\code{rag-dense}) \emph{falla} en tareas causales
entre archivos ($0\%$; cobertura $0{,}19$): la similitud no puede hacer aflorar
contexto causalmente esencial pero léxicamente distinto --- la premisa de la
hipótesis. (Las victorias en $d{=}1$ de \code{rag-hybrid} vienen de la \emph{puntuación
causal} que toma prestada, no de la similitud léxica, razón por la cual persisten bajo
ruido.)
\textbf{(2)} El valor de la selección consciente de la estructura \emph{crece con el
diámetro}: la expansión a un salto solo resuelve $d{=}1$, mientras que la paginación
causal completa (\code{graphrag-bfs}, \code{ccos-region}) resuelve todos los $d$ --- la
dirección de H2.
\textbf{(3) El empate limpio.} Bajo una consulta limpia, \code{graphrag-bfs},
\code{ccos-from-query} y \code{ccos-region} alcanzan todas $1{,}00$: la palanca es la
\emph{estructura} causal, \textbf{no CCOS en sí}.
\textbf{(4) La separación ruidosa.} Bajo una consulta \emph{engañosa}, todo método que
localiza código léxicamente colapsa a $0\%$ --- incluido el fuerte \code{graphrag-bfs}
\emph{y} la ablación \code{ccos-from-query} (CCOS sembrado en la consulta). Solo
\code{ccos-region}, que se ancla en la señal del espacio de trabajo en vez de la
consulta, sobrevive a $1{,}00$. La ablación aísla el diferenciador: es la \emph{fuente
del ancla} (una señal estructural del espacio de trabajo frente a una consulta léxica),
no la maquinaria de regiones. Este es el régimen realista --- una descripción de tarea
rara vez nombra la causa distante exacta, y una memoria de tipo SO que rastrea el
working set activo es robusta donde la recuperación en tiempo de consulta es engañada.
\textbf{(5) Supuestos, declarados.} Esto acredita a CCOS con un ancla fiable (el
archivo activo / la prueba que falla), que una memoria a nivel de SO tiene pero un
\emph{pipeline} de pura recuperación no; y supone repositorios \emph{modulares} con
regiones separables (densidad $0{,}955$ en código real, \S\ref{sec:eval}) --- una
región gigante monolítica que exceda el presupuesto colapsa la ventaja (observado,
reportado).
\textbf{(6)} En todo momento esto es una \emph{simulación bajo un oráculo declarado}:
prueba la condición \emph{necesaria} (recuperación), no la \emph{suficiente}
(generación) de la \S\ref{sec:protocol}.

\section{Evaluación con LLM real: una primera medición en dispositivo}\label{sec:protocol}

La simulación (\S\ref{sec:sim}) estableció la condición necesaria. La condición
suficiente --- que un agente LLM \emph{resuelva} luego más tareas con memoria por
regiones que con recuperación por fragmentos --- requiere despliegues de modelo.
\textbf{Implementamos el arnés} (\code{ccos eval}) y reportamos una \textbf{primera
ejecución con LLM real} contra un modelo local.

\paragraph{Arnés implementado.} Cada tarea es un diminuto proyecto multiarchivo que
codifica una \emph{cadena causal aritmética}: una constante base en un archivo se
transforma a través de una cadena de funciones de una línea en archivos separados, y la
pregunta «¿qué entero devuelve la última función?» solo es respondible \emph{leyendo}
toda la cadena (la causa distante incluida). La calificación es por tanto coincidencia
exacta sobre un entero --- objetiva y automatizable, sin ejecución de código. Las mismas
seis estrategias de la \S\ref{sec:sim} ensamblan la ventana de contexto bajo un
presupuesto de tokens (con la división de consulta limpia/ruidosa), la ventana se envía
a un modelo (cualquier endpoint compatible con OpenAI, Anthropic-Messages u Ollama), y
registramos tres métricas por diámetro: \textbf{éxito de tarea} (entero correcto),
\textbf{cobertura del oráculo} (cadena requerida $\subseteq$ ventana --- independiente
del modelo), y \textbf{alucinación de símbolo} (la respuesta cita una función ausente
del proyecto).

\paragraph{Montaje.} Ejecutamos \code{ccos eval} en dispositivo sobre una NVIDIA Jetson
AGX Thor contra \code{qwen2.5:7b-instruct} servido por Ollama (temperatura $0$), con
$20$ tareas, semilla $7$, un presupuesto de $600$ tokens, diámetros de $1$ a $4$, en
ambos regímenes limpio y ruidoso. La Tabla~\ref{tab:real} reporta la ejecución.

\begin{table}[h]
\centering
\begin{tabular}{lcccccr}
\toprule
& \multicolumn{4}{c}{éxito al diámetro $d$} & & \\
\cmidrule(lr){2-5}
Estrategia & $d{=}1$ & $d{=}2$ & $d{=}3$ & $d{=}4$ & cob. & tok. \\
\midrule
\multicolumn{7}{l}{\emph{Consulta limpia (nombra la función objetivo)}}\\
\code{rag-dense}      & 0.12 & 0.00 & 0.00 & 0.00 & 0.00 & 519 \\
\code{rag-hybrid}     & 0.00 & 0.00 & 0.00 & 0.00 & 0.00 & 521 \\
\code{graphrag-1hop}  & 1.00 & 0.00 & 0.00 & 0.00 & 0.40 & 270 \\
\code{graphrag-bfs}   & 1.00 & 0.00 & 0.17 & 0.00 & 1.00 & 402 \\
\code{ccos-from-query}& 1.00 & 0.00 & 0.17 & 0.00 & 1.00 & 402 \\
\code{ccos-region}    & 1.00 & 0.00 & 0.17 & 0.00 & 1.00 & 402 \\
\midrule
\multicolumn{7}{l}{\emph{Consulta ruidosa (un señuelo supera léxicamente al objetivo)}}\\
\code{rag-dense}      & 0.00 & 0.00 & 0.00 & 0.00 & 0.00 & 531 \\
\code{rag-hybrid}     & 0.00 & 0.00 & 0.00 & 0.00 & 0.00 & 521 \\
\code{graphrag-1hop}  & 0.00 & 0.00 & 0.00 & 0.00 & 0.00 & 165 \\
\code{graphrag-bfs}   & 0.00 & 0.00 & 0.00 & 0.00 & 0.00 & 165 \\
\code{ccos-from-query}& 0.00 & 0.00 & 0.00 & 0.00 & 0.00 & 165 \\
\textbf{\code{ccos-region}} & \textbf{1.00} & \textbf{0.00} & \textbf{0.17} & \textbf{0.00} & \textbf{1.00} & \textbf{402} \\
\bottomrule
\end{tabular}
\caption{Primera ejecución con LLM real: \code{qwen2.5:7b-instruct} vía Ollama, en
dispositivo (NVIDIA Jetson AGX Thor), $20$ tareas, semilla $7$, presupuesto $600$
tokens. «cob.» es la cobertura del oráculo independiente del modelo; «tok.» la media de
tokens de entrada. El éxito es una respuesta entera por coincidencia exacta.}
\label{tab:real}
\end{table}

\paragraph{Hallazgos.}
\textbf{(1) La cobertura se transfiere de la simulación al texto real.} La columna de
cobertura independiente del modelo reproduce la Tabla~\ref{tab:sim} sobre contenido de
archivo real: el RAG léxico cubre $0$, los métodos conscientes de la estructura cubren
$1{,}00$ bajo una consulta limpia, y bajo ruido \emph{solo} \code{ccos-region} mantiene
$1{,}00$ --- la condición necesaria, confirmada fuera del simulador.
\textbf{(2) El éxito está acotado por el modelo, no por el contexto.} Donde la cobertura
es $1{,}00$, el modelo de $7$B aún resuelve solo $d{=}1$ ($1{,}00$) y parte de $d{=}3$
($0{,}17$), fallando $d{=}2,4$: con toda la cadena en la ventana, el fallo residual es
el \emph{razonamiento aritmético} --- un techo de capacidad del modelo, no un fallo de
selección. Crucialmente, ningún método tiene éxito jamás \emph{sin} cobertura (éxito
$\subseteq$ cobertura), que es la afirmación que el arnés está construido para probar.
\textbf{(3) La separación ruidosa se mantiene en un LLM real.} Bajo una consulta
engañosa, todo método impulsado por la consulta colapsa a $0\%$ de éxito --- incluidos
\code{graphrag-bfs} y la ablación \code{ccos-from-query} --- mientras que
\code{ccos-region}, anclada en la señal del espacio de trabajo, es la \emph{única}
estrategia con éxito no nulo ($d{=}1$ a $1{,}00$). Los métodos impulsados por la
consulta incluso parecen \emph{más baratos} bajo ruido ($165$ frente a $402$ tokens)
precisamente porque recuperan con confianza el archivo equivocado, más pequeño;
\code{ccos-region} gasta su presupuesto en la cadena causalmente correcta.
\textbf{(4) Límites honestos.} $20$ tareas ($\approx 5$ por diámetro) es una muestra
pequeña, y un modelo de $7$B pone en el suelo la condición suficiente. Un modelo de
frontera (\code{deepseek-v4-pro}, en curso) o un modelo local de $70$B se espera que
eleve el éxito allí donde la cobertura es $1{,}00$, \emph{agudizando} --- no cambiando
--- la separación, ya que la cobertura fija el techo.

\paragraph{Hacia benchmarks externos.} La suite de cadena aritmética aísla la pregunta
de la selección; las pruebas externas decisivas son la resolución de incidencias de
SWE-bench~\cite{jimenez2023swebench} (un parche pasa las pruebas ocultas) y una suite
controlada de \emph{errores multiarchivo} donde el sitio de la falla, su causa y su
radio de impacto residen en archivos distintos. Las referencias comparten el LLM base y
el presupuesto de tokens: RAG clásico~\cite{lewis2020rag} (top-$k$ coseno sobre
fragmentos), GraphRAG~\cite{edge2024graphrag} (resúmenes de comunidades sobre un grafo
de código), MemGPT~\cite{packer2023memgpt} (memoria paginada de tipo SO), un agente
LangGraph~\cite{langgraph} con un almacén vectorial, y las regiones de CCOS. Métricas:
tasa de resolución, tokens de entrada hasta el éxito, y una tasa de alucinación de
fundamentación de citas verificada contra el grafo verdad-base. Hemos comenzado esto
sobre historia real: \code{scripts/causal\_validation} extrae commits de corrección,
recupera el árbol previo a la corrección, inyecta la falla, y puntúa
$R_{\mathrm{cov}} = |F_{\mathrm{target}} \cap \mathrm{WorkingSet}_K| / |F_{\mathrm{target}}|$
--- la fracción de los archivos que una corrección tocó que el working set acotado
recupera. Sobre la historia de este repositorio, la primera ejecución expuso una
limitación y su corrección, ambas medidas: con propagación de fallo solo aguas abajo,
$R_{\mathrm{cov}}$ era plano en $0{,}33$ (solo el archivo semilla recuperado, ya que los
archivos co-modificados son importadores \emph{aguas arriba} alcanzados solo por
concentradores de dependencias). Resolver las importaciones intra-crate en aristas
archivo$\to$archivo y propagar el fallo \emph{bidireccionalmente} corrige esto. En tres
crates maduras --- \code{fd}, \code{bat} e \code{hyperfine}, $70$ commits de corrección
extraídos --- el efecto es consistente: a presupuesto suficiente ($K{\ge}50$) los dos
cambios elevan $R_{\mathrm{cov}}$ a $0{,}85$--$1{,}0$ (desde $0{,}50$--$0{,}84$ solo
aguas abajo), mientras se \emph{diluye} a $0{,}19$--$0{,}28$ a un presupuesto ajustado
$K{=}20$. \textbf{Pero frente a la referencia obvia, esto no es una victoria.}
Ejecutando un RAG léxico clásico (coseno TF-IDF sobre el texto de los archivos,
consultado por el archivo de la falla) al \emph{mismo presupuesto de archivos},
$R_{\mathrm{cov}}$ para CCOS\,/\,RAG es $0{,}92/0{,}94$, $1{,}00/0{,}98$, $0{,}87/0{,}92$
a $K{=}50$ y empata a $K{=}100$, mientras que a $K{=}20$ el RAG va claramente por
delante ($0{,}20/0{,}56$ en \code{bat}, $0{,}20/0{,}73$ en \code{hyperfine}). La
selección causal no tiene por tanto \emph{ninguna ventaja neta de cobertura} sobre la
similitud léxica aquí, y es peor a presupuesto ajustado: en errores reales los archivos
de una corrección se parecen léxicamente entre sí, así que el TF-IDF también los
recupera --- la premisa de la \S\ref{sec:sim} de una causa léxicamente distinta de su
síntoma no se reproduce en estos repositorios. La lectura honesta es que la condición
\emph{necesaria} se cumple para la gran mayoría pero \emph{no} es específica de CCOS;
una ventaja real tendría que venir de los regímenes que este montaje no prueba --- una
consulta degradada/ausente (\S\ref{sec:sim}), un síntoma (prueba que falla) léxicamente
lejos de su causa (que necesita la siembra real guiada por la prueba, no la heurística
de mayor grado), o la condición \emph{suficiente} (Fase 4). Los espacios de trabajo
multi-crate, ahora enlazables, también quedan por medir a escala.

\paragraph{Condición suficiente (Fase 4): la resolución empata, la eficiencia no.}
Ejecutamos la mitad de generación sobre correcciones reales de un solo archivo de
\code{fd}: a un agente se le da un contexto de igual presupuesto construido de dos
formas --- la región causal de CCOS frente a un RAG léxico top-$k$ --- se le pide
reescribir el archivo con el error, y se califica según si \code{cargo test} pasa, con
un reintento \emph{compilador en el bucle} (al fallar, una \emph{falla de página de
contexto} analiza el error con la capa de traza, inyecta presión en los archivos que
fallan, y vuelve a indicar con una ventana refrescada). Con un modelo débil de $7$B y un
solo intento, ambos resuelven $0/15$; con \code{qwen3-coder-30B} y tres intentos el
bucle eleva ambos a $2/10$ ($20\%$) --- es la retroalimentación por falla de página, no
el recuperador, lo que desbloquea la resolución, y CCOS no muestra \textbf{ninguna
ventaja de resolución} sobre el RAG, consistente con el resultado de recuperación. La
\emph{eficiencia} los separa: a igual resolución, la región auto-acotada de CCOS
promedia $776$ tokens de contexto frente a los $5366$ que llenan el presupuesto del
recuperador léxico --- una reducción de $\mathbf{6{,}9\times}$. Lo mismo se cumple a
escala sin modelo: en $51$ escenarios de corrección de un solo archivo de \code{fd},
\code{bat} e \code{hyperfine}, CCOS ensambla $700$--$1600$ tokens de contexto frente a
los $\approx\!6000$ del RAG, una reducción de $\mathbf{4{,}1}$--$\mathbf{9{,}1\times}$.
Este es el único eje en el que CCOS domina: no \emph{lo que} recupera (una referencia lo
iguala) sino \emph{cuán poco} necesita, porque la región causal se detiene en el working
set en lugar de rellenar un top-$k$ hasta el presupuesto. Enunciamos las salvedades: la
muestra de \emph{resolución} es diminuta ($n{=}10$, dos pasadas), y la referencia llena
el presupuesto por construcción (un RAG con $k$ cuidadosamente ajustado también sería
más disperso). La afirmación defendible es más estrecha pero real: CCOS
\emph{se autocalibra} --- se acota en la región causal sin $k$ ni presupuesto que
ajustar --- así que nunca desperdicia la ventana, que es precisamente el propósito de la
paginación bajo demanda.

\paragraph{Hipótesis.}
\textbf{H1} (eficiencia): CCOS alcanza un éxito de tarea igual con menos tokens de
entrada que RAG/GraphRAG, porque una región cubre el vecindario causal con mayor
exhaustividad (Tabla~\ref{tab:locality}).
\textbf{H2} (éxito a largo horizonte): en tareas multiarchivo, la tasa de resolución de
CCOS supera al RAG por fragmentos por un margen que crece con el diámetro causal de la
tarea.
\textbf{H3} (fundamentación): CCOS reduce la tasa de alucinación de símbolos, porque la
ventana admitida es un subgrafo conexo de nodos \emph{reales}.

\paragraph{Amenazas a la validez.} La calidad de una región está acotada por la calidad
de las aristas (\S\ref{sec:eval}); los resultados confundirían la \emph{política} de
región con el \emph{grafo} subyacente. Las ablaciones deben, por tanto, fijar el grafo y
variar solo la estrategia de selección, y reportar el diámetro causal por tarea para que
las ganancias se atribuyan a la coherencia, no a recuperar más texto. La ablación
\code{ccos-from-query} de la Tabla~\ref{tab:real} es exactamente este control --- mismo
grafo y presupuesto, solo el ancla intercambiada --- y colapsa con las referencias
léxicas bajo ruido. Un resultado positivo en H1--H3 constituiría la contribución de
investigación; un resultado nulo aún validaría la infraestructura determinista y
auditable.

\section{Limitaciones}

Somos deliberadamente explícitos sobre lo que \emph{no} se muestra.
\textbf{(1) Ninguna ventaja de cobertura sobre RAG en errores reales.} La métrica
principal sobre datos reales ($R_{\mathrm{cov}}$, \S\ref{sec:protocol}) empata con un
simple recuperador léxico TF-IDF a presupuesto suficiente y pierde ante él a uno
ajustado; en commits de corrección reales los archivos que una corrección toca se
parecen léxicamente, así que la premisa de la simulación de una causa léxicamente
distinta no se sostiene. Los fuertes números absolutos son una condición
\emph{necesaria} que una referencia también satisface, no una victoria de CCOS.
(Una medición complementaria sobre la propia fuente del motor aísla una relación
distinta, definida estructuralmente: un recuperador léxico TF-IDF recupera solo el
$49\%$ (exhaustividad@10) de las dependencias de importación resueltas por AST ---
pares que rutinariamente no comparten vocabulario, coseno de dependencia $0{,}53$
frente a $0{,}43$ aleatorio --- así que el punto ciego del RAG es real, pero queda
fuera del eje de cohesión de la corrección que mide la métrica principal, donde los
archivos de una corrección sí comparten vocabulario. La capa estructural recupera
esas aristas por construcción; el AST, ahora por defecto, es lo que hace completa esa
recuperación.)
\textbf{(2) Necesario $\neq$ suficiente.} Nunca medimos si una región \emph{ayuda a un
agente a corregir} un error (Fase 4: un parche que pasa las pruebas ocultas) --- solo si
los archivos relevantes son recuperables.
\textbf{(3) Heurística de siembra.} El arnés inyecta la falla en el archivo modificado
de mayor grado de salida, no en la prueba realmente fallida, así que no ejercita el
régimen donde un síntoma está léxicamente lejos de su causa --- el régimen más favorable
a la selección causal.
\textbf{(4) Escala y estadística.} Tres repositorios de una sola crate,
$n\!\approx\!20$ cada uno; las diferencias se reportan con su desviación estándar pero no
se validan de forma cruzada.
\textbf{(5) Motor.} El analizador ahora usa por defecto un AST \code{syn} real ---
medido un $36{,}5\%$ más preciso que la antigua heurística de línea sobre la propia
fuente del motor (dos tercios frente a exhaustividad de importación completa), con la
heurística conservada solo como reserva para entradas que no son Rust --- y la
ingesta se endureció para reconstruir de forma bit-idéntica entre procesos
(resolución de importaciones ordenada; centralidad invariante al orden). Las regiones
siguen siendo de granularidad de archivo y solo se fusionan en aristas entre archivos
explícitas; los pesos de puntuación están ajustados a mano, aunque ya están
disponibles refinamientos, desactivados por defecto, de centralidad de autovalor y de
ciclo de vida de nodo (Stable/Working/Orphan). Nada de esto afecta a las propiedades
\emph{probadas} (partición, determinismo, reproducción) ni a la localidad medida ---
pero el caso empírico de un beneficio para el agente aguas abajo está, en este punto,
\emph{sin probar}.

\section{Conclusión}

Nos propusimos mostrar que organizar la memoria de un agente por \emph{regiones
causales} recupera el contexto a largo horizonte mejor que la recuperación por
fragmentos, dimos a la construcción una definición precisa y determinista, y probamos
que se reconstruye bit a bit en la reproducción. Luego probamos la afirmación de
recuperación contra la referencia obvia sobre datos reales --- y no se sostuvo: en $70$
commits reales de corrección de errores, un simple recuperador léxico TF-IDF empata con
la selección causal y la supera a presupuesto ajustado, y un pivote por traza de fallo
pierde ante un RAG-sobre-el-mensaje-de-error, porque en código real los archivos de una
corrección y sus mensajes de error comparten vocabulario. Reportamos el resultado
negativo en vez de enterrarlo. Lo que sobrevive no es un mejor recuperador sino un tipo
de objeto distinto: una memoria de trabajo \emph{determinista, reproducible, auditable}
en la que el estado exacto del contexto de un agente puede rebobinarse y reproducirse
bajo parámetros distintos --- depuración por viaje en el tiempo que una pila de
recuperación probabilística no puede ofrecer. Si esa auditabilidad, o una ventaja de
contexto estructurado en tiempo de \emph{generación} (Fase 4), produce una ganancia
medible aguas abajo, queda abierto. Publicamos el motor, el arnés de validación honesto
y este artículo para que la distinción entre lo que está probado, lo que está medido y
lo que está refutado pueda comprobarse en vez de afirmarse.

\begin{thebibliography}{99}
\bibitem{lewis2020rag} P.~Lewis et al. Retrieval-Augmented Generation for
Knowledge-Intensive NLP Tasks. \emph{NeurIPS}, 2020. arXiv:2005.11401.
\bibitem{asai2023selfrag} A.~Asai et al. Self-RAG: Learning to Retrieve, Generate
and Critique through Self-Reflection. \emph{ICLR}, 2024. arXiv:2310.11511.
\bibitem{edge2024graphrag} D.~Edge et al. From Local to Global: A Graph RAG
Approach to Query-Focused Summarization. 2024. arXiv:2404.16130.
\bibitem{yamaguchi2014cpg} F.~Yamaguchi et al. Modeling and Discovering
Vulnerabilities with Code Property Graphs. \emph{IEEE S\&P}, 2014.
\bibitem{packer2023memgpt} C.~Packer et al. MemGPT: Towards LLMs as Operating
Systems. 2023. arXiv:2310.08560.
\bibitem{park2023generative} J.~S.~Park et al. Generative Agents: Interactive
Simulacra of Human Behavior. \emph{UIST}, 2023. arXiv:2304.03442.
\bibitem{langgraph} LangChain. LangGraph: Building Stateful, Multi-Actor
Applications with LLMs. Marco de software, 2023--2024.
\url{https://github.com/langchain-ai/langgraph}.
\bibitem{denning1968working} P.~J.~Denning. The Working Set Model for Program
Behavior. \emph{Communications of the ACM}, 11(5), 1968.
\bibitem{liu2023lost} N.~F.~Liu et al. Lost in the Middle: How Language Models Use
Long Contexts. \emph{TACL}, 2024. arXiv:2307.03172.
\bibitem{kwon2023paged} W.~Kwon et al. Efficient Memory Management for Large
Language Model Serving with PagedAttention. \emph{SOSP}, 2023. arXiv:2309.06180.
\bibitem{haber1991timestamp} S.~Haber and W.~S.~Stornetta. How to Time-Stamp a
Digital Document. \emph{Journal of Cryptology}, 3(2), 1991.
\bibitem{jimenez2023swebench} C.~E.~Jimenez et al. SWE-bench: Can Language Models
Resolve Real-World GitHub Issues? \emph{ICLR}, 2024. arXiv:2310.06770.
\end{thebibliography}

\end{document}
